% Capítulo 5 – gerado a partir do relatório; edite livremente
\chapter{PARÂMETROS DE CORTE}

\section{Equações utilizadas}

Os parâmetros foram calculados com as relações da Aula 2. A rotação é obtida da velocidade de corte e do diâmetro, e a velocidade de avanço, do avanço por volta ou por dente:


\begin{equation}\label{eq:n}
  n = \frac{1000 \cdot v_{c}}{\pi \cdot d}
\end{equation}

\begin{equation}\label{eq:vf}
  v_{f} = f \cdot n
\end{equation}

\begin{equation}\label{eq:vfz}
  v_{f} = f_{z} \cdot z \cdot n
\end{equation}

em que n é a rotação (rpm), v\textsubscript{c} a velocidade de corte (m/min), d o diâmetro da peça ou da ferramenta (mm), f o avanço por volta (mm), f\textsubscript{z} o avanço por dente (mm) e z o número de dentes da ferramenta. No acabamento, a rugosidade máxima teórica depende do avanço e do raio de ponta r\textsubscript{$\varepsilon$} (Aula 10):


\begin{equation}\label{eq:rmax}
  R_{\text{máx}} = \frac{f^{2}}{8 \cdot r_{\varepsilon}}
\end{equation}

No fresamento com pequena penetração lateral, o cavaco fica mais fino que o avanço por dente. A espessura máxima do cavaco foi calculada com o ângulo de contato $\varphi$\textsubscript{0}:


\begin{equation}\label{eq:phi0}
  \cos \varphi_{0} = 1 - \frac{2 \cdot a_{e}}{D}
\end{equation}

\begin{equation}\label{eq:hex}
  h_{ex} = f_{z} \cdot \sin \varphi_{0}
\end{equation}

\section{Critério de escolha}

Os valores de partida foram tirados das recomendações dos fabricantes para o grupo P2.1. No torno, porém, a potência e o torque disponíveis limitam esses valores. O Centur 30D com cabeçote de 3000 rpm tem motor de 9 kW, mas só fornece essa potência acima de cerca de 600 rpm; abaixo disso o limite é o torque de 140 N$\cdot$m. Como a peça tem Ø160 mm, o desbaste ocorre entre 550 e 600 rpm, justamente na faixa limitada pelo torque. A Tabela~\ref{tab:verif} mostra que os valores nominais de catálogo excediam a máquina. Por isso o avanço de desbaste foi reduzido para 0,22 mm/volta, o faceamento foi dividido em dois passes e, na furação e no rebaixo, foram adotados os menores valores recomendados pelo fabricante. O cálculo das potências está no Capítulo 6.


\begin{table}[H]
\centering
\caption{Potência exigida com os valores nominais de catálogo no torno Centur 30D}\label{tab:verif}
{\small\renewcommand{\arraystretch}{1.2}
\begin{tabular}{|>{\raggedright\arraybackslash}p{\dimexpr 0.2709\linewidth-2\tabcolsep\relax}|>{\centering\arraybackslash}p{\dimexpr 0.0985\linewidth-2\tabcolsep\relax}|>{\centering\arraybackslash}p{\dimexpr 0.0862\linewidth-2\tabcolsep\relax}|>{\centering\arraybackslash}p{\dimexpr 0.1108\linewidth-2\tabcolsep\relax}|>{\centering\arraybackslash}p{\dimexpr 0.0985\linewidth-2\tabcolsep\relax}|>{\centering\arraybackslash}p{\dimexpr 0.1601\linewidth-2\tabcolsep\relax}|>{\centering\arraybackslash}p{\dimexpr 0.1601\linewidth-2\tabcolsep\relax}|}
\hline
\rowcolor{cinza} \textbf{Operação} & \textbf{a\textsubscript{p} (mm)} & \textbf{f (mm)} & \textbf{v\textsubscript{c} (m/min)} & \textbf{n (rpm)} & \textbf{Potência exigida (kW)} & \textbf{Potência disponível (kW)} \\ \hline
Faceamento de desbaste & 2,85 & 0,30 & 335 & 666 & 22,3 & 9,0 \\ \hline
Torneamento externo de desbaste & 1,50 & 0,30 & 335 & 666 & 11,7 & 9,0 \\ \hline
Furação Ø23 & 11,50 & 0,13 & 200 & 2.768 & 14,5 & 9,0 \\ \hline
Rebaixo frontal & 5,00 & 0,12 & 180 & 546 & 10,7 & 8,0 \\ \hline
\end{tabular}}
\fonte{Elaborado pelos autores.}
\end{table}


\section{Parâmetros adotados}

A Tabela~\ref{tab:param} apresenta os parâmetros adotados nas operações com ferramenta de geometria definida. Como exemplo, no faceamento de desbaste, partindo do diâmetro de 160 mm com v\textsubscript{c} = 300 m/min, as Equações \eqref{eq:n} e \eqref{eq:vf} dão:


\[
  n = \frac{1000 \cdot 300}{\pi \cdot 160} = 597 \ \text{rpm}
\]

\[
  v_{f} = 0,22 \cdot 597 = 131 \ \text{mm/min}
\]

\begin{table}[H]
\centering
\caption{Parâmetros de corte adotados}\label{tab:param}
{\small\renewcommand{\arraystretch}{1.2}
\begin{tabular}{|>{\raggedright\arraybackslash}p{\dimexpr 0.3264\linewidth-2\tabcolsep\relax}|>{\centering\arraybackslash}p{\dimexpr 0.0699\linewidth-2\tabcolsep\relax}|>{\centering\arraybackslash}p{\dimexpr 0.1282\linewidth-2\tabcolsep\relax}|>{\centering\arraybackslash}p{\dimexpr 0.1166\linewidth-2\tabcolsep\relax}|>{\centering\arraybackslash}p{\dimexpr 0.1049\linewidth-2\tabcolsep\relax}|>{\centering\arraybackslash}p{\dimexpr 0.1282\linewidth-2\tabcolsep\relax}|>{\centering\arraybackslash}p{\dimexpr 0.1107\linewidth-2\tabcolsep\relax}|}
\hline
\rowcolor{cinza} \textbf{Operação} & \textbf{Ferr.} & \textbf{a\textsubscript{p} (mm)} & \textbf{f (mm)} & \textbf{v\textsubscript{c} (m/min)} & \textbf{n (rpm)} & \textbf{v\textsubscript{f} (mm/min)} \\ \hline
20.1 / 30.1 – Faceamento de desbaste (2 passes) & T01 & 1,425 & 0,22 & 300 & 597 a 3.000 & 131 \\ \hline
20.2 / 30.3 – Torneamento externo de desbaste & T01 & 1,5 & 0,22 & 280 & 557 & 122 \\ \hline
20.3 / 30.2 – Degrau do aro & T01 & 1,15 & 0,22 & 300 & 608 a 2.032 & – \\ \hline
20.4 / 30.4 – Acabamento & T02 & 0,5 & 0,15 & 400 & 811 a 3.000 & 122 \\ \hline
20.5 – Furação Ø23 & T03 & 11,5 & 0,09 & 150 & 2.076 & 187 \\ \hline
20.6 – Mandrilamento Ø24,5 & T04 & 0,75 & 0,10 & 231 & 3.000 & 300 \\ \hline
20.7 / 30.5 – Rebaixo frontal & T05 & 5 (largura) & 0,07 & 160 & 485 a 999 & – \\ \hline
40 – Fresamento dos dentes & F01 & 12,6 e 0,4 (a\textsubscript{e}) & 0,055 (f\textsubscript{z}) & 18,8 & 60 & 39,6 \\ \hline
50 – Rasgos: desbaste & F02 & 5 & 0,10 (f\textsubscript{z}) & 275 & 5.471 & 1.094 \\ \hline
50 – Rasgos: acabamento & F02 & 5 (a\textsubscript{e} = 1) & 0,20 (f\textsubscript{z}) & 275 & 5.471 & 2.188 \\ \hline
\end{tabular}}
\fonte{Elaborado pelos autores com base nos catálogos Sandvik Coromant e ZPS-FN.}
\end{table}


No mandrilamento, a velocidade recomendada, de 395 a 465 m/min, exigiria mais de 5.000 rpm; com o limite de 3.000 rpm do cabeçote, a velocidade fica em 231 m/min, o que apenas aumenta um pouco o tempo de corte. No fresamento dos dentes, os 15 m/min recomendados corresponderiam a 48 rpm, abaixo da rotação mínima da fresadora, de 60 rpm. Com 60 rpm a velocidade é de 18,8 m/min, valor entre os indicados pelo fabricante para aços de até 850 N/mm² (15 m/min) e de até 600 N/mm² (20 m/min). No acabamento dos rasgos, com 1 mm de penetração lateral, o avanço por dente foi aumentado para 0,20 mm, de modo que a espessura máxima do cavaco (Equação \eqref{eq:hex}) fique em 0,097 mm, próxima da recomendada.


Nas retificações, os parâmetros estão na Tabela~\ref{tab:paramret}. Na retificação plana, o rebolo de 350 mm a 1.450 rpm tem velocidade periférica de 26,6 m/s; na retificação interna, o rebolo de Ø20 mm precisa girar a cerca de 28.648 rpm para atingir 30 m/s, valor dentro da faixa dos eixos internos da S121.


\begin{table}[H]
\centering
\caption{Parâmetros das operações de retificação}\label{tab:paramret}
{\small\renewcommand{\arraystretch}{1.2}
\begin{tabular}{|>{\raggedright\arraybackslash}p{\dimexpr 0.1947\linewidth-2\tabcolsep\relax}|>{\centering\arraybackslash}p{\dimexpr 0.0916\linewidth-2\tabcolsep\relax}|>{\raggedright\arraybackslash}p{\dimexpr 0.1833\linewidth-2\tabcolsep\relax}|>{\raggedright\arraybackslash}p{\dimexpr 0.1833\linewidth-2\tabcolsep\relax}|>{\centering\arraybackslash}p{\dimexpr 0.1833\linewidth-2\tabcolsep\relax}|>{\centering\arraybackslash}p{\dimexpr 0.1489\linewidth-2\tabcolsep\relax}|}
\hline
\rowcolor{cinza} \textbf{Operação} & \textbf{v\textsubscript{s} (m/s)} & \textbf{Peça} & \textbf{Avanço} & \textbf{Penetração por passe (mm)} & \textbf{Passes} \\ \hline
90 – Retificação plana & 26,6 & v\textsubscript{w} = 15 m/min & 8 mm por curso & 0,02 / 0,005 & 7 + 2 + 2 sem penetração \\ \hline
100 – Retificação interna & 30 & v\textsubscript{w} = 30 m/min (382 rpm) & 10 e 5 mm por volta & 0,005 / 0,0025 & 44 + 12 + 4 sem penetração \\ \hline
\end{tabular}}
\fonte{Elaborado pelos autores.}
\end{table}

